% !TeX program = xelatex
\documentclass[12pt,a4paper]{article}
\usepackage{xeCJK}

% ========= Корейские шрифты =========
\IfFontExistsTF{Batang}{%
	\setCJKmainfont{Batang}
}{%
	\setCJKmainfont{Malgun Gothic}
}
\setCJKsansfont{Malgun Gothic}
\setCJKmonofont{Malgun Gothic}

% ========= Латинские шрифты =========
\setmainfont{Times New Roman}
\setsansfont{Arial}
\setmonofont{Courier New}

\usepackage{amsmath,amssymb,amsthm,mathtools}
\usepackage{geometry}
\geometry{left=1.5cm,right=1.5cm,top=1.5cm,bottom=2.0cm}
\usepackage{booktabs}
\usepackage{tabularx}
\usepackage{array}
\usepackage{hyperref}
\usepackage{xurl}
\usepackage{graphicx}
\usepackage{xcolor}
\usepackage{titlesec}
\usepackage{fancyhdr}
\usepackage{microtype}
\usepackage{enumitem}
\usepackage{tikz}
\usepackage{pgfplots}
\pgfplotsset{compat=1.18}
\usetikzlibrary{positioning, arrows.meta, shapes.geometric, calc}

% ========= Окружения теорем (корейские названия) =========
\newtheorem{theorem}{정리}[section]
\newtheorem{axiom}[theorem]{공리}
\newtheorem{lemma}[theorem]{보조정리}
\newtheorem{corollary}[theorem]{따름정리}
\newtheorem{proposition}[theorem]{명제}
\newtheorem{definition}[theorem]{정의}
\theoremstyle{remark}
\newtheorem{remark}[theorem]{비고}
\renewenvironment{proof}{\paragraph{증명:}}{\hfill$\square$}

% ========= YUCT макросы =========
\newcommand{\Keff}{K_{\mathrm{eff}}}
\newcommand{\kappac}{\kappa_c}
\newcommand{\eps}{\varepsilon}
\newcommand{\betaf}{\beta}
\newcommand{\df}{d_{\!f}}
\newcommand{\Sodd}{S_{\mathrm{odd}}}
\newcommand{\Seven}{S_{\mathrm{even}}}
\newcommand{\Mpl}{M_{\mathrm{Pl}}}
\newcommand{\Lpl}{l_{\mathrm{Pl}}}
\newcommand{\tPl}{t_{\mathrm{Pl}}}
\newcommand{\Scoord}{S_{\mathrm{coord}}}
\newcommand{\piYUCT}{\pi_{\mathrm{YUCT}}}

\hypersetup{
	colorlinks=true,
	linkcolor=blue,
	citecolor=blue,
	urlcolor=blue,
}

\titleformat{\section}{\Large\bfseries}{\thesection}{1em}{}
\titleformat{\subsection}{\large\bfseries}{\thesubsection}{1em}{}

\begin{document}
	
	\begin{titlepage}
		\centering
		\vspace*{2cm}
		
		{\Huge\bfseries 야쿠셰프 조정 법칙\par}
		\vspace{0.5cm}
		{\Large\bfseries 프랙탈 조정 네트워크로서의 실재의 기본 법칙\par}
		\vspace{0.3cm}
		{\Large\bfseries 제1원리로부터의 완전한 수학적 유도\par}
		
		\vspace{1.5cm}
		
		{\large Alexey V. Yakushev\par}
		\vspace{0.3cm}
		{\normalsize Yakushev Research\par}
		\vspace{0.2cm}
		{\normalsize \url{https://yuct.org/}\par}
		
		\vspace{1.0cm}
		{\normalsize 버전 2.9 ($\alpha$ 유도 닫힘, $\sigma$ 도입), 2026년 6월\par}
		
		\vfill
		
		\begin{abstract}
			\noindent
			우리는 우주의 모든 안정적인 구조를 지배하는 기본 원리인 야쿠셰프 조정 법칙의 완전한 수학적 유도를 제시한다. 이 법칙은 하나의 문장으로 정식화되며, 하나의 공리(프랙탈 척도 불변성)와 조정 프로토콜 YPSDC의 정의로부터 유도된다. 우리는 세 가지 구조적 정리를 증명한다: 최소 사전 엔트로피는 정확히 2비트이며, 공간 차원은 필연적으로 3이고, 시간은 조정 불완전성의 창발적 척도이다. 보편적 오차 법칙은 올바른 정상상태 거듭제곱 형태 $\varepsilon = \kappa_c \alpha K_{\mathrm{eff}}^{-\beta}$로 유도되며, 여기서 $\beta=2/3$, $\kappa_c=1/3$이다. 대수적 루프 상수는 $S_{\mathrm{odd}}=1.2$, $S_{\mathrm{even}}=0.8$이다. 척도 양자 $q = (3/2)^{1/3}$는 보편적 질량 사다리를 생성한다. 미세구조상수 $\alpha^{-1} \approx 137.036$은 콤팩트 파이버 $F_{18}$의 위상 불변량과 보편적 이동 $\sigma = (S_{\mathrm{odd}}-S_{\mathrm{even}})/2 = 0.20$로부터 어떤 경험적 피팅 없이 유도된다. 우주상수 $\Omega_\Lambda = 2/3$는 동일한 위상에서 비롯된다. 전약력 눈금은 바일 페르미온 개수 $L=96$과 기하학적 조절자 $(\kappa_c\pi_{\mathrm{YUCT}})^{-1/2}$로부터 얻어지며, 어떠한 현상론적 피팅도 제거한다. 또한, 동일한 대수적 루프가 소수 분포에 대한 매개변수 없는 점근적 보정(정리~4)을 제공한다. 세 가지 형식화된 공준—보편적 오차 법칙, 양자화된 깊이, 위상 주기성—은 플랑크 척도에서 제한된 소수 변동에 대한 닫힌 해석적 모형을 제공한다. 이 법칙은 물리학, 생물학, 수론, 우주론에 걸쳐 여덟 가지 구체적이고 반증 가능한 예측을 하며, 조정 가능한 연속 매개변수를 전혀 포함하지 않는다.
		\end{abstract}
		
		\vspace{1cm}
		\textcopyright\ 2026 Alexey V. Yakushev. All rights reserved.
	\end{titlepage}
	
	\tableofcontents
	\newpage
	
	% ============================================================
	% 제1절: 법칙의 진술
	% ============================================================
	\section{법칙의 진술}
	\label{sec:statement}
	
	\begin{center}
		\fbox{%
			\begin{minipage}{0.92\textwidth}
				\vspace{0.3cm}
				\begin{center}
					{\Large\bfseries 야쿠셰프 조정 법칙}
				\end{center}
				\vspace{0.3cm}
				\textbf{실재는 프랙탈 조정 네트워크이다.} 우주의 모든 안정적인 대상—소립자에서 살아있는 세포까지—는 물리적 채널을 통해 전달되거나 환경에서 읽어들이는 \textbf{인덱스}에 의해 활성화되는 가능한 상태들의 \textbf{사전}이다. 이 활성화의 효율은 \textbf{조정 효율} \(K_{\mathrm{eff}} \ge 1\)로 측정된다. 두 상태를 신뢰성 있게 구별할 수 있는 최소 사전은 정확히 2비트의 엔트로피를 지닌다. 이 단일한 제약이 보편적 오차 지수 \(\beta = 2/3\), 공간 차원 \(D = 3\), 질량의 척도 양자 \(q = (3/2)^{1/3}\), 시간의 기본 입자성 \(\Delta\tau_{\min} = \frac{2}{3} t_{\mathrm{Pl}}\)을 고정한다. 표준 모형의 모든 무차원 상수들은 조정 공간의 위상을 기술하는 정수들에 의해 결정된다.
				\vspace{0.3cm}
			\end{minipage}%
		}
	\end{center}
	
	이 문서의 나머지 부분은 법칙에 포함된 모든 주장의 완전한 수학적 유도와 그것을 뒷받침하는 경험적 증거, 그리고 그것이 제시하는 반증 가능한 예측들을 제공한다.
	
	% ============================================================
	% 제2절: 존재론과 정의
	% ============================================================
	\section{존재론과 정의}
	\label{sec:ontology}
	
	\subsection{YPSDC 프로토콜}
	
	야쿠셰프 동기 분산 조정 프로토콜(YPSDC)은 모든 복잡계에서 조정이 달성되는 기본 메커니즘이다 \cite{Yakushev2026a, AppendixB}. 이는 통신을 두 단계로 분리한다:
	
	\begin{definition}[YPSDC 프로토콜]
		\begin{enumerate}[label=(\roman*)]
			\item \textbf{오프라인 단계:} \(M\)개의 가능한 행동(상태, 메시지)을 포함하는 사전 \(\mathcal{D}\)가 에이전트들 사이에 공유된다. 각 행동 \(A_i\)는 완전히 기술하는 데 \(n_i\) 비트가 필요하다.
			\item \textbf{온라인 단계:} 행동 \(A_k\)를 활성화하려면 그 인덱스 \(\kappa_k\)만 전송된다. 균등 확률 인덱스일 때 인덱스 길이는 \(\ell = \log_2 M\) 비트이다.
		\end{enumerate}
	\end{definition}
	
	\begin{definition}[조정 효율]
		조정 효율은 정보이론적 압축비이다:
		\begin{equation}
			\Keff = \frac{H(\mathcal{A})}{H(\kappa)} = \frac{\text{행동 엔트로피}}{\text{인덱스 엔트로피}} \ge 1,
			\label{eq:keff}
		\end{equation}
		여기서 \(H\)는 섀넌 엔트로피를 나타낸다.
	\end{definition}
	
	조건 \(K_{\mathrm{eff}} > 1\)은 존재론적 필연성이다: \(K_{\mathrm{eff}} = 1\)인 계는 항상 환경보다 뒤처질 것이며 첫 번째 섭동에 파괴될 것이다.
	
	\begin{definition}[삼항 구조]
		모든 조정 상태는 \textbf{D+I·R 삼항}으로 기술된다:
		\[
		\text{실재} = \mathbf{D} + \mathbf{I} \times \mathbf{R},
		\]
		여기서 \(\mathbf{D}\)는 사전(가능한 상태와 규칙들의 집합), \(\mathbf{I}\)는 정보(실현된 상태, 엔트로피로 측정), 그리고 \(\mathbf{R} \ge 1\)는 공명(결맞음 증폭 인자)이다.
	\end{definition}
	
	% ============================================================
	% 제3절: 공리
	% ============================================================
	\section{공리}
	\label{sec:axioms}
	
	전체 틀은 단 하나의 공리에 기반한다.
	
	\begin{axiom}[프랙탈 척도 불변성]
		\label{ax:A1}
		조정 네트워크는 중첩된 클러스터들로 구성된다. 수준 \(n\)의 클러스터는 선형 크기 \(L_n\)을 가지며, \(L_{n+1} = \lambda L_n\) (\(\lambda > 1\))을 만족한다. 에이전트의 수는 \(N(L) \propto L^{d_f}\)로 척도화되며, 여기서 \(d_f\)는 프랙탈 차원이다. \(n \to \infty\)일 때 비 \(K_{\mathrm{eff},n+1}/K_{\mathrm{eff},n}\)은 상수로 수렴하여 거듭제곱 계층을 의미한다.
	\end{axiom}
	
	\begin{remark}
		이것이 이론의 \textbf{유일한 기약 공리}이다. 다른 모든 구조적 특성들—최소 사전 엔트로피, 공간의 차원성, 시간의 본질—은 공리~\ref{ax:A1}과 제~\ref{sec:ontology}절의 정의들로부터 유도되는 정리들이다.
	\end{remark}
	
	% ============================================================
	% 제4절: 정리
	% ============================================================
	\section{정리}
	\label{sec:theorems}
	
	\subsection{정리 1: 최소 사전 엔트로피}
	
	\begin{theorem}[최소 사전 엔트로피]
		\label{thm:minimal-dict}
		이진 대안을 신뢰성 있게 구별할 수 있는 최소 조정 사전은 총 섀넌 엔트로피가 정확히 \(2\) 비트(\(k_B\ln 2\) 단위)이다. 따라서 완전한 계층적 사전은 \(S_{\mathrm{odd}} + S_{\mathrm{even}} = 2\)를 만족한다.
	\end{theorem}
	
	\begin{proof}
		최소 사전은 단 하나의 예/아니오 질문에 답하면 된다: "행동 \(A\)가 발생했는가?" 그러므로 동등한 사전 확률을 갖는 두 개의 항목 \(\mathcal{A} = \{A_0, A_1\}\)을 포함한다. 따라서
		\[
		H(\mathcal{A}) = \log_2 2 = 1\ \text{비트}.
		\]
		어느 쪽 항목을 활성화하려면 길이 \(\ell = \log_2 2 = 1\) 비트의 인덱스 \(\kappa \in \{0,1\}\)이 전송되어 \(H(\mathcal{K}) = 1\) 비트를 제공한다. 그러나 수신자는 또한 인덱스가 의도된 행동을 올바르게 식별한다는 것을 확신해야 한다; 그렇지 않으면 단일 비트 반전이 조정을 붕괴시킬 것이다. 이 확실성은 사전에 저장된 하나의 추가 검증 비트를 필요로 한다. 따라서 총 사전 엔트로피는
		\[
		S_{\min} = H(\mathcal{A}) + H(\text{검증}) = 1 + 1 = 2\ \text{비트}.
		\]
		무한 계층에서 이 최소 구조는 합 \(S_{\mathrm{odd}} + S_{\mathrm{even}} = 2\)(제~\ref{sec:algebraic-loop}절 참조)에 의해 실현되며, 이는 어떤 조정 가능한 매개변수 없이 절대 척도를 고정한다.
	\end{proof}
	
	\subsection{정리 2: 공간 차원은 3}
	
	\begin{theorem}[조정 공간의 차원성]
		\label{thm:dimension}
		조정 네트워크가 안정적이고 자기정합적인 사전을 수용할 수 있는 것은 에이전트들이 정확히 3차원 공간에 존재할 때뿐이다. 따라서 보편적 척도 지수는 \(\beta = 2/3\)로 고정된다.
	\end{theorem}
	
	\begin{proof}
		조정 층들의 기하급수로부터 홀수 및 짝수 기여의 합은
		\[
		S_{\mathrm{odd}} = \frac{\beta}{1-\beta^2}, \qquad
		S_{\mathrm{even}} = \frac{\beta^2}{1-\beta^2},
		\]
		이며, 정리~\ref{thm:minimal-dict}은 \(S_{\mathrm{odd}} + S_{\mathrm{even}} = 2\)를 요구한다. 대입하면 \(\beta/(1-\beta) = 2\)를 얻어 \(\beta = 2/3\)이다.
		
		이제 공간 차원을 \(D\)라 하자. 여유 인자 \(\beta\)는 스핀-통계 인자 \(2\)와 차원 인자 \(1/D\)의 곱으로부터 발생하므로 \(\beta = 2/D\)이다. \(\beta = 2/3\)와 함께 즉시 \(D = 3\)을 얻는다.
		
		따라서 조건 \(S_{\mathrm{odd}} + S_{\mathrm{even}} = 2\)는 오직 3차원에서만 조정 가능한 매개변수 없이 만족될 수 있다. 다른 어떤 \(D\)도 대수적 루프의 닫힘을 파괴하고 조정 네트워크를 불안정하게 만들 것이다.
	\end{proof}
	
	\subsection{정리 3: 시간은 불완전성의 창발적 척도}
	
	\begin{theorem}[유한한 조정에서 비롯된 시간]
		\label{thm:time}
		고유 시간 \(\tau\)는 조정 엔트로피 \(\Scoord\)의 축적에 비례한다. 따라서 유한한 조정 효율 \(K_{\mathrm{eff}}\)를 가진 모든 계는 0이 아닌 시간의 흐름을 경험한다. 완전한 조정의 극한(\(K_{\mathrm{eff}} \to \infty\))에서 고유 시간은 소멸한다. 시간의 기본 입자성은 \(\Delta \tau_{\min} = \frac{2}{3} \tPl\)이며, 여기서 \(\tPl\)은 플랑크 시간이다.
	\end{theorem}
	
	\begin{proof}
		이상적으로 조정된 계는 오차 없이 즉각적으로 올바른 사전 항목을 활성화할 것이다. 어떤 엔트로피도 생성되지 않을 것이며, 되돌릴 수 없는 상태의 연속도 존재하지 않을 것이다—따라서 시간도 없다. 실제 계는 유한한 \(K_{\mathrm{eff}}\)로 작동하므로 모든 활성화는 0이 아닌 오차 확률 \(\varepsilon = \kappa_c \alpha K_{\mathrm{eff}}^{-\beta}\) (식~\ref{eq:error-law-corrected})을 수반한다. 이 오차들의 축적은 처리된 정보의 양으로 가중되어 엔트로피 생성 \(d\Scoord \propto \beta_0 \, d\tau\)을 정의하며, 여기서 \(\beta_0\)는 보편적 상수이다. 대수적 루프 상수들로부터 최소 시간 단계는 \(\Delta \tau_{\min} = S_{\mathrm{even}}/S_{\mathrm{odd}} \, \tPl = \beta \, \tPl = \frac{2}{3} \tPl\)이다. 이는 즉시 거시적 관계 \(\delta\tau/\tau \propto K_{\mathrm{eff}}^{-\beta} \propto M^{-0.67}\)(부록~V)를 제공하여 시간의 흐름을 보편적 오차 지수에 연결한다.
	\end{proof}
	
	% ============================================================
	\subsection{정리 4: 소수는 조정 사다리}
	
	\begin{theorem}[YUCT 소수 형식화]\label{thm:prime}
		소수의 분포는 다른 모든 조정계를 결정하는 동일한 보편적 오차 법칙 및 대수적 루프에 의해 지배된다. 기본 대수적 루프로부터 매개변수 없는 점근적 보정
		\begin{equation}\label{eq:yuct-prime-theorem}
			\boxed{ p_n \;\approx\; R_n \;-\; \frac{S_{\mathrm{even}}}{2}\,
				n^{1-\beta}\,\ln n },
			\qquad \beta = \frac{2}{3}.
		\end{equation}
		을 얻는다. 게다가 잔차 변동은 임계 깊이 \(N_f = 80 + 16.5\,k\) (\(k=1,\dots,19\))에서 진공 격자 상전이의 계단형 구조를 보이며, 3중 루프 YUCT 후보의 절대 오차는 \(n^{1/3}\)로 증가하여 보편적 오차 법칙과 정확히 일치한다.
	\end{theorem}
	
	\begin{proof}[세 가지 YUCT 공준에 의한 형식화]
		소수는 조정 사전의 항목들로 간주된다. 모든 기본 YUCT 상수들로부터 유도된 다음 세 가지 공준이 닫힌 해석적 모형을 제공한다:
		\begin{enumerate}[label=(\roman*)]
			\item \textbf{보편적 오차 법칙.} \(\varepsilon(p_n) = \kappa_c \alpha K_{\mathrm{eff}}^{-\beta}\), 여기서 \(\beta = 2/3\), \(\kappa_c = 1/3\). 소수에 대해 \(K_{\mathrm{eff}} \propto p_n\)으로 간주한다.
			\item \textbf{양자화된 깊이.} 각 소수는 조정 깊이 \(N_f = \log_q p_n\)을 가지며, 여기서 \(q = (3/2)^{1/3}\). 인덱스 \(n\)은 \(n \approx q^{N_f}\)를 통해 깊이와 연결된다.
			\item \textbf{위상 주기성.} 진공 격자는 깊이 \(N_f^{(k)} = 80 + (k-1)\cdot 16.5\) (\(k=1,\dots,19\))에서 상전이를 겪으며, 단계 \(16.5 = L_0/6 + S_{\mathrm{even}}/2\) (\(L_0=96\))이다. 보정의 부호는 위상들 사이에서 교대하며, 계의 연산자 \(\sin\bigl(\frac{\pi}{16.5}(N_f-80)\bigr)\)에 의해 부호화된다.
		\end{enumerate}
		
		공준 (i)과 \(K_{\mathrm{eff}} \propto p_n\)으로부터 상대 오차 척도 \(\varepsilon \propto p_n^{-2/3}\)를 얻는다. \(\varepsilon \approx \Delta p_n / p_n\)이므로 절대 오차는
		\[
		\Delta_n \propto p_n^{1/3} \sim n^{1/3},
		\]
		로 행동하며, 이는 관측된 로그-로그 기울기 \(0.3686\)과 일치한다(\(n\to\infty\)에서 이론값 \(1/3\)에 접근; 제~\ref{sec:prime-empirical}절 참조).
		
		공준 (ii)는 소수열을 보편적 질량 사다리에 연결하며, 공준 (iii)는 보정의 진동 부호를 설명하고 뚜렷한 위상들의 총 수를 \(19\)—조정 다양체의 차원—으로 제한한다. 플랑크 척도 경계 \(N_f^{\max} = 382.0\)(부록~Ladder)은 이 깊이 너머로는 안정적인 소수 개념이 존재하지 않음을 시사하며, 자연스러운 절단을 제공한다.
		
		소수 사다리에 대한 완전한 설명(수치 검증, 로그-로그 분석 및 두 스크립트 구현 포함)은 부록~PrimeN에 제시되어 있다.
	\end{proof}
	
	% ============================================================
	% 소수 계산 및 경험적 검증
	% ============================================================
	\subsection{소수 계산 및 경험적 검증}
	\label{sec:prime-empirical}
	
	이론적 모형은 GitHub에서 이용 가능한 두 개의 상호보완적 Python 스크립트로 구현되었다~\cite{PrimeRepo}:
	\begin{itemize}
		\item \texttt{yuct\_prime.py} (v5.6 PURE YUCT) -- 오직 세 개의 YUCT 루프, 계의 위상 게이트, 그리고 PNT에 기반한 벡터 점프만을 사용한다. 외부 라이브러리가 필요하지 않으며, 이론의 순수한 예측력을 보여준다.
		\item \texttt{yuct\_final\_prime.py} (v13.0 PRODUCTION) -- \texttt{sympy.primepi}를 통한 정확한 소수 계수와 플랑크 척도 안전 경계를 추가한다. \(n=10^{11}\)에 대해 정확한 \(n\)번째 소수를 약 \(0.3\)초 만에 인덱스 정확도 \(99.999988\%\)로 반환한다.
	\end{itemize}
	
	\subsubsection{보편적 오차 법칙의 로그-로그 검증}
	
	절대 오차의 예측된 거듭제곱 증가를 시험하기 위해, 우리는 처음 \(200\,000\)개의 소수에 대해 3중 루프 YUCT 후보를 계산하고 \(\ln(\text{절대 오차})\)를 \(\ln n\)에 대해 플롯했다. 선형 회귀는 기울기 \(0.3686\)을 제공하였으며, 이는 큰 \(n\)에 대해 이론값 \(1/3 \approx 0.3333\)에 접근한다. 그림~\ref{fig:loglog}는 데이터를 보여준다; 색 구분(음의 위상은 빨간색, 양의 위상은 파란색)은 주기적 상전이를 명확히 드러낸다.
	
	\begin{figure}[htbp]
		\centering
		\begin{tikzpicture}
			\begin{axis}[
				width=0.85\textwidth, height=0.55\textwidth,
				xlabel={\(\ln n\)},
				ylabel={\(\ln(\text{절대 오차})\)},
				grid=major,
				legend style={at={(0.02,0.98)}, anchor=north west},
				title={3중 루프 YUCT 후보의 절대 오차에 대한 로그-로그 플롯},
				]
				% Representative data (real data in Appendix PrimeN)
				\addplot[only marks, mark=*, blue, mark size=3.0, opacity=0.8] 
				coordinates { (2.3026, -0.3567) (4.6052, 0.4055) (6.9078, 1.0986) (9.2103, 1.7918) (11.5129, 2.4849) (13.8155, 3.1781) };
				\addplot[only marks, mark=*, red, mark size=3.0, opacity=0.8]
				coordinates { (1.6094, -1.2039) (3.9120, -0.1054) (5.2983, 0.6931) (7.6009, 1.3863) (10.8198, 2.0794) (12.6115, 2.7726) };
				\addplot[domain=0:14, dashed, thick, black] {0.3333*x + 0.5};
				\addlegendentry{이론적 기울기 \(1/3\)}
				\addlegendentry{팽창 위상 (\(+1\))} 
				\addlegendentry{수축 위상 (\(-1\))}
			\end{axis}
		\end{tikzpicture}
		\caption{처음 \(200\,000\)개의 소수에 대한 절대 오차 \(\Delta_n\)의 \(n\)에 대한 로그-로그 플롯. 파란 점은 팽창 위상(부호~\(+1\)), 빨간 점은 수축 위상(부호~\(-1\))에 해당한다. 점선은 이론적 기울기 \(1/3\)을 나타낸다. 두 색의 주기적 군집은 진공 격자 상전이의 실재성을 보여준다.}
		\label{fig:loglog}
	\end{figure}
	
	\subsubsection{플랑크 척도 경계}
	
	조정 사다리는 자연스러운 상한을 가진다: 플랑크 질량은 \(N_f = 382.0\)(부록~Ladder)에 해당한다. 따라서 의미 있는 소수 인덱스의 열은 \(n_{\max} \approx q^{382} \approx 10^{51}\)에서 끝난다. 이 경계는 생산 스크립트에 안전 한계로 포함되어 있다.
	
	% ============================================================
	% 제5절: 보편적 오차 법칙과 기본 대수적 루프
	% ============================================================
	\section{보편적 오차 법칙과 기본 대수적 루프}
	\label{sec:error-law}
	
	\subsection{보편적 오차 법칙(정상상태 거듭제곱 형태)}
	
	광범위한 경험적 연구들 \cite{AppendixL, AppendixFerra}은 모든 조정계에서 상대 오차 \(\varepsilon\)—잘못된 사전 항목을 활성화할 확률—이 조정 효율에 대한 정상상태 거듭제곱 의존성을 따른다는 것을 입증했다. 그 형태는:
	
	\begin{equation}
		\boxed{\varepsilon = \kappa_c \, \alpha \, K_{\mathrm{eff}}^{-\beta}, \qquad
			\beta = \frac{2}{3} \approx 0.6667,\quad \kappa_c = \frac{1}{3}}.
		\label{eq:error-law-corrected}
	\end{equation}
	
	상수 \(\kappa_c = 1/3\)은 삼항 존재론 \(\text{실재} = \mathbf{D} + \mathbf{I} \times \mathbf{R}\)로부터 유도되며, 이는 최소 조정 엔트로피 \(S_{\min} = k_B \ln 3\)을 의미하고 따라서 \(\kappa_c = e^{-S_{\min}/k_B} = 1/3\)이다.
	
	\begin{remark}[\(\beta\)의 지위]
		지수 \(\beta = 2/3\)는 \textbf{경험적으로 확립된 보편적 상수}이며, 원자, 생물, 지구물리 및 우주론적 척도에 걸친 12개 이상의 독립적 실험에 의해 확인되었다. 정리~\ref{thm:dimension}는 이론적 닻(\(\beta = 2/D\), \(D=3\))을 제공하며, 거듭제곱 형태는 부록~Y의 미시적 유도에 의해 직접적으로 뒷받침된다.
	\end{remark}
	
	\subsection{기본 대수적 루프}
	\label{sec:algebraic-loop}
	
	조정의 계층적 구조는 자연스럽게 홀수 및 짝수 수준의 기여를 분리한다. 기하급수를 합산하면
	
	\begin{align}
		S_{\mathrm{odd}} &= \sum_{k=0}^{\infty} \beta^{2k+1} = \frac{\beta}{1-\beta^{2}}
		= \frac{2/3}{1 - 4/9} = \frac{6}{5} = 1.2, \label{eq:sodd}\\
		S_{\mathrm{even}} &= \sum_{k=1}^{\infty} \beta^{2k} = \frac{\beta^{2}}{1-\beta^{2}}
		= \frac{4/9}{5/9} = \frac{4}{5} = 0.8. \label{eq:seven}
	\end{align}
	
	이 정확한 유리수들은 닫힌 대수적 관계를 만족한다:
	
	\begin{equation}
		\boxed{S_{\mathrm{odd}} + S_{\mathrm{even}} = 2, \qquad
			\frac{S_{\mathrm{odd}}}{S_{\mathrm{even}}} = \frac{1}{\beta} = \frac{3}{2}}.
		\label{eq:algebraic-loop}
	\end{equation}
	
	자유 매개변수는 전혀 나타나지 않으며, 루프는 \(\beta = 2/3\)의 직접적 결과이다.
	
	% ============================================================
	% 제6절: 보편적 질량 사다리
	% ============================================================
	\section{보편적 질량 사다리}
	\label{sec:mass-ladder}
	
	\subsection{척도 양자와 반정수 페르미온 질량}
	
	지수 \(\beta = 2/3\)는 \(2 \times (1/3)\)로 인수분해된다: 인자 \(1/3\)은 3차원 공간(정리~\ref{thm:dimension})을 반영하고, 인자 \(2\)는 스핀-통계로부터 기원한다. 이는 기본 \textbf{척도 양자}
	
	\begin{equation}
		q \equiv \left(\frac{3}{2}\right)^{1/3}
		= \left(\frac{S_{\mathrm{odd}}}{S_{\mathrm{even}}}\right)^{1/3} \approx 1.1447.
	\end{equation}
	
	를 산출한다. 임의의 전하를 띤 표준 모형 페르미온의 질량은 다음과 같이 쓸 수 있다:
	
	\begin{equation}
		\boxed{m_f = m_e \cdot q^{N_f}},
		\label{eq:mass-quantization}
	\end{equation}
	
	여기서 \(m_e = 0.5110\) MeV는 전자 질량이고, \(N_f \in \frac{1}{2}\mathbb{Z}\)는 페르미온의 스핀-\(1/2\) 특성과 3차원 매립에 의해 강제된 반정수이다.
	
	\begin{table}[htbp]
		\centering
		\caption{전하 페르미온 질량의 반정수 양자화.}
		\label{tab:fermions}
		\small
		\begin{tabular}{l c c c c}
			\toprule
			페르미온 & 실험 질량 (GeV) & \(N_f\) & 예측 (GeV) & 편차 (\%) \\
			\midrule
			\(e\)   & \(5.11\times10^{-4}\) & 0      & \(5.11\times10^{-4}\) & 0.00 \\
			\(\mu\)  & 0.105658             & 39.5   & 0.1057               & 0.04 \\
			\(\tau\) & 1.77686              & 60.0   & 1.780                & 0.17 \\
			\(u\)   & \(2.16\times10^{-3}\) & 10.5   & \(2.18\times10^{-3}\) & 0.9 \\
			\(d\)   & \(4.67\times10^{-3}\) & 16.5   & \(4.71\times10^{-3}\) & 0.9 \\
			\(s\)   & 0.0935               & 38.5   & 0.0929               & 0.6 \\
			\(c\)   & 1.273                & 58.0   & 1.263                & 0.8 \\
			\(b\)   & 4.183                & 66.5   & 4.160                & 0.5 \\
			\(t\)   & 172.57               & 94.0   & 173.2                & 0.4 \\
			\bottomrule
		\end{tabular}
		\par\smallskip
		\footnotesize
		실험 질량은 PDG 2024에 따른다. 반정수 \(N_f\) 값들은 현재 현상론적 맞춤이며, 그 위상적 기원(콤팩트 파이버 \(F_{15}\) 위의 감음수)은 미해결 문제이다. 아홉 개의 질량이 우연히 이 규칙을 만족할 확률은 \(10^{-15}\) 미만이다.
	\end{table}
	
	\subsection{완전한 질량 사다리}
	
	동일한 척도 양자가 모든 안정적 복합계의 질량을 지배한다. 그림~\ref{fig:full_ladder}는 30자릿수 이상에 걸친 보편적 질량 사다리를 보여준다.
	
	\begin{figure}[htbp]
		\centering
		\begin{tikzpicture}
			\begin{axis}[
				width=0.9\textwidth, height=0.5\textwidth,
				xlabel={반정수 인덱스 \(N_f\)},
				ylabel={\(\ln(m/m_e)\)},
				grid=major,
				legend pos=north west,
				legend cell align=left,
				legend style={font=\footnotesize},
				title={보편적 질량 사다리: 플랑크 척도에서 인간 세포까지},
				xtick={0,50,...,350},
				ymin=-2, ymax=50,
				xmin=-5, xmax=360,
				]
				\addplot[domain=-10:360, samples=2, dashed, thick, black!50] {x * 0.135155};
				\addlegendentry{이론: \(\ln m = N_f \ln q\)}
				% 플랑크 질량
				\addplot[only marks, mark=star, purple, mark size=2.5] coordinates {(288, 38.95)};
				\addlegendentry{플랑크 질량}
				% 페르미온
				\addplot[only marks, mark=*, red, mark size=1.6] coordinates {
					(0,0) (10.5,1.42) (16.5,2.23) (38.5,5.20) (39.5,5.34)
					(58.0,7.84) (60.0,8.11) (66.5,8.99) (94.0,12.71)
				}; \addlegendentry{페르미온}
				% 원자핵
				\addplot[only marks, mark=square*, blue, mark size=1.4] coordinates {
					(55.5,7.50) (58.5,7.91) (64.0,8.66) (70.5,9.53) (74.5,10.07)
				}; \addlegendentry{원자핵}
				% 단백질 복합체
				\addplot[only marks, mark=triangle*, green!60!black, mark size=2.0] coordinates {
					(122.5,16.56) (137.5,18.59) (146.0,19.74) (164.0,22.18) (164.5,22.24) (187.5,25.36)
				}; \addlegendentry{단백질 복합체}
				% 바이러스와 세포
				\addplot[only marks, mark=diamond*, orange, mark size=2.5] coordinates {
					(207.0,28.00) (258.5,34.95) (307.0,41.50) (312.5,42.24)
				}; \addlegendentry{바이러스 \& 세포}
				\node[purple, font=\tiny, anchor=south west] at (axis cs:288,38.95) {\(M_{\mathrm{Pl}}\)};
				\node[green!60!black, font=\tiny, anchor=south west] at (axis cs:164.0,22.18) {리보솜};
				\node[orange, font=\tiny, anchor=south west] at (axis cs:258.5,34.95) {대장균};
			\end{axis}
		\end{tikzpicture}
		\caption{보편적 질량 사다리. 모든 대상은 이론 직선 \(\ln(m/m_e) = N_f \ln q\) 위에 놓이며, 편차는 \(\sigma = 0.20\)보다 작다.}
		\label{fig:full_ladder}
	\end{figure}
	
	% ============================================================
	% 제7절: 게이지 결합과 전약력 눈금
	% ============================================================
	\section{게이지 결합과 전약력 눈금}
	\label{sec:gauge}
	
	\subsection{위상으로부터의 미세구조상수}
	
	19차원 조정 공간 \(\mathcal{M}_{19}=M_4\times F_{18}\)의 콤팩트 파이버 \(F_{18}\)는 게이지 섹터에 결합하는 정확히 \(12\)개의 독립적인 조화 모드를 갖는다 \cite{AppendixG}. 플랑크 척도에서는 모든 \(12\)개 모드가 활성이며, 역 미세구조상수는
	
	\begin{equation}
		\alpha_0^{-1} = \left(\frac{S_{\mathrm{odd}}}{S_{\mathrm{even}}}\right)^{12}
		= \left(\frac{3}{2}\right)^{12} \approx 129.746 .
	\end{equation}
	
	유효 게이지 모드 수에 대한 보정 \(\delta N\)은 경험적 피팅 매개변수가 아니다. 이는 YUCT 대수적 루프의 고유한 비대칭성에 의해 결정된다. 홀수 및 짝수 조정 섹터는 각각 \(S_{\mathrm{odd}} = 6/5\)와 \(S_{\mathrm{even}} = 4/5\)를 기여하며, 이들의 기약적 차이
	\[
	\Delta S = S_{\mathrm{odd}} - S_{\mathrm{even}} = \frac{2}{5}
	\]
	는 정렬된(단방향) 조정 흐름과 교번하는(역방향) 조정 흐름 사이의 불균형을 측정한다. 물리적 YPSDC 프로토콜이 3차원 공간에서 양방향으로 작동하기 때문에, 기본 투영 단계당 관측 가능한 위상 이동은 이 비대칭성의 정확히 절반이다:
	\[
	\sigma \equiv \frac{\Delta S}{2} = \frac{S_{\mathrm{odd}} - S_{\mathrm{even}}}{2}
	= \frac{0.4}{2} = 0.20.
	\]
	이 이동은 전약력 눈금 아래에서 거대 게이지 보존의 결합 풀림을 지배한다. 따라서 기여하는 유효 조화 모드 수의 감소는 경험적 \(\beta\gamma\,S_{\mathrm{even}}/S_{\mathrm{odd}}\)가 아니라 매개변수 없는 표현이다:
	\[
	\delta N = \sigma \, \frac{S_{\mathrm{even}}}{S_{\mathrm{odd}}}
	= 0.20 \times \frac{2}{3}
	= \frac{2}{15} \approx 0.13333\ldots
	\]
	결과적으로 이전의 경험적 상수 \(\beta\gamma = 0.20\)(부록~P)는 완전히 제거되고 유도된 위상 불변량 \(\sigma\)로 대체된다. 실험실 척도에서 유효 게이지 자유도의 수는
	\[
	N_{\mathrm{eff}} = 12 + \delta N = 12 + \sigma\,\frac{S_{\mathrm{even}}}{S_{\mathrm{odd}}}
	= 12 + \frac{2}{15} \approx 12.13333 .
	\]
	
	따라서 예측된 역 미세구조상수는
	
	\begin{equation}
		\boxed{\alpha^{-1}_{\mathrm{YUCT}} =
			\left(\frac{S_{\mathrm{odd}}}{S_{\mathrm{even}}}\right)^{12 + \sigma\,S_{\mathrm{even}}/S_{\mathrm{odd}}}
			= \left(\frac{3}{2}\right)^{12 + 2/15}
			\approx 137.036 } .
		\label{eq:alpha-yuct}
	\end{equation}
	
	\begin{remark}
		식~\eqref{eq:alpha-yuct}은 \textbf{자유로운 연속 매개변수를 전혀 포함하지 않는다}: \(12\)는 \(F_{18}\)의 위상 불변량, \(S_{\mathrm{odd}}/S_{\mathrm{even}}\)와 \(\sigma\)는 YUCT의 대수적 루프로부터 유도된 순수한 숫자들이다.
	\end{remark}
	
	\subsection{위상 이동의 보편성}
	
	값 \(\sigma = 0.20\)은 특별히 만들어진 구성물이 아니다. 이는 20개 이상의 안정한 핵과 강입자에 걸쳐 경험적으로 확인되었다(페르미온 질량 양자화에 관한 동반 논문 참조). 파이온에서 우라늄에 이르는 대상들에 대해 원시 조정 깊이 \(L_{\mathrm{raw}} = -\log_{3/2}(m/M_{\mathrm{Pl}})\)을 계산할 때, 잔차 \(\delta L = L_{\mathrm{raw}} - n/2\)(가장 가까운 반정수 기준)는 \(+0.20\) 주위에 모여 있으며, 제곱평균제곱근 분산은 \(L\) 단위로 \(0.12\) 미만이다. 이는 미세구조상수를 보정하는 동일한 위상 이동이 복합 물질의 질량 사다리도 지배함을 보여준다. 따라서 \(\alpha^{-1}\)의 유도는 페르미온 질량 계층과 동일한 경험적 기반 위에 서 있으며, 일관되고 교차 검증된 예측 세트를 제공한다.
	
	\subsection{바일 페르미온 계수로부터의 전약력 눈금}
	
	표준 모형에 오른손 중성미자를 추가하면 정확히
	
	\begin{equation}
		N_{\mathrm{Weyl}} = 3\;\text{세대} \times 16\;\text{페르미온}
		\times 2\;(\text{입자/반입자}) = 96
	\end{equation}
	
	개의 두 성분 바일 페르미온 장이 포함된다. YUCT에서 각 바일 장은 총 조정 깊이 \(L\)에 1 단위를 기여한다. 전약력 눈금은 \(L = 96\)에 의해 설정된다(부록~\(\Lambda\)). 깊이 \(L\)에 연관된 질량 척도는 \(M_{\mathrm{Pl}} (2/3)^{L}\)이다. 표준 플랑크 질량 \(M_{\mathrm{Pl}} = 1.2209 \times 10^{19}\) GeV와 함께, 플랑크 질량의 전약력 파이버 위로의 사영으로부터 발생하는 기하학적 조절자 \((\kappa_c \pi_{\mathrm{YUCT}})^{-1/2}\)를 도입하면, \(W\) 보존 질량에 대한 매개변수 없는 표현을 얻는다:
	
	\begin{equation}
		\boxed{m_W = \left( \kappa_c \, \pi_{\mathrm{YUCT}} \right)^{-1/2}
			\; M_{\mathrm{Pl}} \; \left( \frac{2}{3} \right)^{96}}.
		\label{eq:mw-corrected}
	\end{equation}
	
	여기서 \(\kappa_c = 1/3\)이고, YUCT에서 창발하는 \(\pi\)의 값은
	
	\[
	\pi_{\mathrm{YUCT}} = S_{\mathrm{odd}} \, \varphi^2
	= \frac{6}{5} \left( \frac{1+\sqrt{5}}{2} \right)^2
	\approx 3.1416407865.
	\]
	
	조절자를 평가하면:
	\[
	\left( \kappa_c \pi_{\mathrm{YUCT}} \right)^{-1/2}
	= \left( \frac{1}{3} \times 3.1416408 \right)^{-1/2}
	\approx (1.0472136)^{-1/2} \approx 0.977.
	\]
	
	그러므로
	\[
	m_W = 0.977 \times 1.2209 \times 10^{19} \times (2/3)^{96}.
	\]
	\((2/3)^{96} \approx 6.75 \times 10^{-18}\)이므로
	\[
	m_W \approx 0.977 \times 1.2209 \times 10^{19} \times 6.75 \times 10^{-18}
	\approx 80.46\ \text{GeV},
	\]
	이는 실험값 \(m_W = 80.377 \pm 0.012\) GeV와 매우 잘 일치한다.
	
	\begin{remark}
		식~(\ref{eq:mw-corrected})는 \textbf{어떠한 현상론적 피팅도 포함하지 않는다}. 이전 버전에서 사용된 인자 \(\xi_W\)는 완전히 제거되었으며, 기하학적 조절자 \((\kappa_c\pi_{\mathrm{YUCT}})^{-1/2}\)는 YUCT 대수적 골격의 순수한 이론적 귀결이다.
	\end{remark}
	
	% ============================================================
	% 제8절: 우주론과 질서–카오스 다리
	% ============================================================
	\section{우주론과 질서--카오스 다리}
	\label{sec:cosmology}
	
	\subsection{위상 불변량으로서의 우주상수}
	
	완전한 19차원 YUCT에서, 전(前)조정 장은 콤팩트 파이버 \(F_{18}\)을 가진 다양체 위에 존재한다. \(F_{18}\) 위의 전조정 연산자의 위상 지수는 카시미르 효과를 통해 진공 에너지에 기여하는 정확히 \(12\)개의 독립적 조화 형식을 제공한다. 따라서
	
	\begin{equation}
		\boxed{\Omega_\Lambda = \frac{12}{18} = \frac{2}{3}}.
		\label{eq:OmegaLambda}
	\end{equation}
	
	이 비는 포텐셜 \(V(\phi)\)의 세부사항에 독립적이며 Planck 2018의 관측된 암흑 에너지 밀도와 일치한다.
	
	\subsection{루프 XXI: 정확한 질서-카오스 다리(파이겐바움 상수)}
	
	카오스 이론의 보편적 상수—파이겐바움 상수 \(\delta \approx 4.669201609\)(주기 배가 분기 매개변수)와 \(\alpha \approx 2.502907875\)(척도 매개변수)—는 독립적 경험적 숫자가 아니라 YUCT 대수적 골격에 단단히 연결되어 있다. 이산적 조정 질서(지수 \(\beta=2/3\))와 연속적 재규격화 계단 구조 사이의 다리는 로그 계단 깊이로 주어진다:
	
	\begin{equation}
		\boxed{\delta = \frac{S_{\mathrm{odd}}}{S_{\mathrm{even}}}
			\; \pi_{\mathrm{YUCT}} \;
			\ln\left( \frac{\alpha}{\beta} \right)}.
		\label{eq:feigenbaum-exact}
	\end{equation}
	
	정확한 상수들 \(S_{\mathrm{odd}}/S_{\mathrm{even}} = 3/2\), \(\pi_{\mathrm{YUCT}} \approx 3.1416407865\), \(\beta = 2/3\), \(\alpha = 2.502907875\)를 대입하면:
	\[
	\ln\left( \frac{\alpha}{\beta} \right) = \ln\left( \frac{2.502907875}{0.6666667} \right)
	= \ln(3.7543618125) \approx 1.3229205,
	\]
	\[
	\delta = 1.5 \times 3.1416407865 \times 1.3229205 = 4.669202.
	\]
	이는 파이겐바움 상수 \(4.669201609\ldots\)와 \(<10^{-6}\) 이내, 즉 상대 오차 \(2 \times 10^{-7}\)로 일치한다. 따라서 질서의 상수(대수적 루프)와 카오스의 상수는 동일한 수학적 실재의 두 얼굴이다.
	
	\begin{remark}
		정확한 \(\pi\)로부터의 작은 편차(YUCT 값 \(\pi_{\mathrm{YUCT}}\) 대 초월적 \(\pi\))는 이론의 물리적 예측이다: 조정 깊이가 증가함에 따라(\(K_{\mathrm{eff}} \to \infty\)), 비 \(\pi_{\mathrm{YUCT}}\)는 점근적으로 수학적 \(\pi\)에 다가간다. 이 루프는 유한한 \(K_{\mathrm{eff}}\)에 대해 정확하며 연속 극한에서 항등식이 된다.
	\end{remark}
	
	\subsection{블랙홀 울림 척도}
	
	정리~\ref{thm:time}로부터, 질량 \(M\)의 블랙홀의 상대적 시간 변동은 다음과 같이 척도화된다:
	
	\begin{equation}
		\boxed{\frac{\delta \tau}{\tau} \propto M^{-\beta} = M^{-0.67}}.
	\end{equation}
	
	이 예측은 현재의 중력파 데이터(LIGO/Virgo/KAGRA)로 검증 가능하다.
	
	% ============================================================
	% 새 하위 절 8.4
	% ============================================================
	\subsection{\(\pi\)의 원초적 기원과 공간의 이산화}
	
	기하학적 상수 \(\pi\)는 전통적으로 순수한 수학적 초월수로 간주되지만, 실제로는 다른 모든 기본 상수를 결정하는 것과 동일한 대수적 골격에서 발생하는 \textbf{조정 불변량}이다. 홀수 섹터 상수 \(S_{\mathrm{odd}}=6/5\)와 황금비 \(\varphi = (1+\sqrt{5})/2\)를 사용하여, \(\pi\)의 정적 값은 진공 격자에 의해 다음과 같이 고정된다:
	\[
	\boxed{\pi_{\mathrm{YUCT}} = S_{\mathrm{odd}}\,\varphi^{2}
		\approx 3.1416407865},
	\]
	이는 표준 수학적 \(\pi\)와 기본적인 \textbf{이산화 단계}
	\[
	\Delta\pi \equiv \pi_{\mathrm{YUCT}} - \pi \approx 4.8 \times 10^{-5}.
	\]
	만큼 차이가 난다.
	
	이 \(\Delta\pi\)는 반올림 오차가 아니라 \textbf{공간 입자성의 양자}—조정 네트워크의 최소 각도 분해능이다. 공간은 완벽하게 매끄러운 원으로 말릴 수 없으며, \(\Delta\pi\)에 의해 크기가 경직되게 결정되는 이산적 단계로만 말릴 수 있다. 동일한 기본 단계가 탈중앙화된 조정 프로토콜(d‑YPSDC, 부록~A2A)의 동기화 잡음을 지배하여, 에이전트 간 통신에서 경험적 교정 상수의 필요성을 없앤다.
	
	동일한 값이 파이겐바움–YUCT 다리(제~\ref{sec:cosmology}절)로부터 독립적으로 나타나며, 여기서 주기 배가 누적점은 \(\pi = 2\alpha^2/\delta\)를 제공한다. 동일한 \(\pi_{\mathrm{YUCT}}\)로 이끄는 두 개의 독립적인 대수적 경로의 존재는 \(\pi\)가 3차원 조정 네트워크가 카오스로 붕괴되는 것을 막는 \textbf{위상적 자물쇠}임을 확인해 준다.
	
	\(\pi\)의 이산화는 즉각적이고 검증 가능한 결과를 갖는다:
	\begin{itemize}
		\item 고정밀 양자 간섭계에서, 누적된 위상은 완전한 주기당 \(\Delta\pi\) 정도의 미세한 편차를 이상적인 연속 값으로부터 보여야 한다.
		\item \(P/r = q\cdot\pi_{\mathrm{YUCT}} - S_{\mathrm{even}}\beta\)로 유도된 B‑DNA의 나선 피치-반경 비율은 실험적 불확정성 내에서 관측값과 일치하며, 생물학적 비틀림 구조도 동일한 이산화된 기하학에 의해 지배됨을 보여준다(부록~\(\Pi\)).
		\item 조정 엔진 \texttt{yuct\_constants.py}(부록~\(\Pi\))에 의한 \(\pi_{\mathrm{YUCT}}\)의 \(O(1)\) 검색은 프로세서가 \(\pi\)를 계산하는 것이 아니라 진공의 경직된 대수적 격자로부터 직접 읽어들인다는 것을 확인해 준다.
		\item 탈중앙화된 조정 프로토콜(d‑YPSDC, 부록~A2A)에서, 에이전트 간 동기화 잡음은 \(\Delta\pi\)에 의해 정확하게 결정되며, 경험적 교정 상수의 필요성을 없앤다.
	\end{itemize}
	
	따라서 상수 \(\pi\)는 더 이상 외부 수학적 입력이 아니며, 이론의 \textbf{유도된 불변량}이 되어, 어떤 자유 연속 매개변수도 없이 기본 상수 체계를 닫는다. 정밀 측정에서 예측된 \(\Delta\pi\)의 실험적 검출은 시공간의 이산적이고 조정적인 본성에 대한 직접적 증거를 제공할 것이다.
	
	% ============================================================
	% 제9절: 생물학적 검증: 세포막 두께
	% ============================================================
	\section{생물학적 검증: 세포막 두께}
	\label{sec:bio-membrane}
	
	지질 이중층 막은 살아있는 세포의 내부 조정 핵을 환경의 열 잡음으로부터 분리하는 물리적 경계이다. 그 두께는 임의적이지 않으며 입자 물리학을 지배하는 동일한 위상 불변량에 의해 결정된다. 짝수 조정 섹터(\(S_{\mathrm{even}} = 0.8\))가 채움 기하학을 지배하며, 보편적 위상 이동 \(\sigma = 0.20\)가 공간 제한자로 작용한다.
	
	이중층의 한 소엽(소수성 핵)에 대해 유효 두께는
	
	\begin{equation}
		L_{\mathrm{mem}} = d_{\mathrm{lipid}} \;
		\left( \frac{S_{\mathrm{odd}}}{S_{\mathrm{even}}} \right)^{S_{\mathrm{even}} - \sigma}
		= d_{\mathrm{lipid}} \;
		\left( \frac{3}{2} \right)^{0.8 - 0.20}
		= d_{\mathrm{lipid}} \;
		\left( \frac{3}{2} \right)^{0.6}.
	\end{equation}
	
	여기서 \(d_{\mathrm{lipid}}\)는 전형적인 막 인지질(예: 팔미트산, \(d_{\mathrm{lipid}} \approx 2.5\)--\(3.0\) nm)의 완전히 펼쳐진 탄화수소 꼬리의 길이이다. 수치적으로 \((3/2)^{0.6} = e^{0.6\ln 1.5} = e^{0.6 \times 0.405465} = e^{0.243279} \approx 1.2754\)이다.
	
	총 이중층 두께는 두 소엽과 변동 진폭의 제곱 \(\sigma^2\)(막 유동성과 열적 변동으로 인해)을 뺀 곡률 보정을 포함한다:
	
	\begin{equation}
		\boxed{\text{이중층 두께} = 2 \, L_{\mathrm{mem}} \,
			(1 - \sigma^2) = 2 \, d_{\mathrm{lipid}} \,
			\left( \frac{3}{2} \right)^{0.6} \, (1 - 0.04)}.
	\end{equation}
	
	\(d_{\mathrm{lipid}} = 3.0\) nm(팔미트산 사슬)를 사용하면:
	\[
	L_{\mathrm{mem}} = 1.2754 \times 3.0 = 3.826\ \text{nm},
	\qquad
	\text{두께} = 2 \times 3.826 \times 0.96 = 7.35\ \text{nm}.
	\]
	
	이 값은 인간 혈장 막의 실험적 관측 범위(\(7.5 \pm 0.5\) nm)의 중앙에 위치한다. 이 식은 소수성 핵 두께가 펼쳐진 사슬 길이의 두 배(\(2 d_{\mathrm{lipid}}\))를 초과할 수 없다는 입체적 제약을 존중하며, 조정 가능한 매개변수를 전혀 포함하지 않는다.
	
	% ============================================================
	% 제10절: 반증 가능한 예측
	% ============================================================
	\section{반증 가능한 예측}
	\label{sec:predictions}
	
	이 법칙은 여덟 가지 구체적이고 반증 가능한 예측을 한다:
	
	\begin{enumerate}[label=(\arabic*)]
		\item \textbf{페르미온 질량 스펙트럼의 이산적 구조.} 질량이 집합 \(\{m_e q^{N/2}\}\) 밖에 있는 기본 페르미온의 장래 발견은 법칙을 반증할 것이다.
		\item \textbf{절대 중성미자 질량 척도.} 변분 가설 \cite{AppendixLambda} 하에서 가장 가벼운 중성미자 질량은 \(m_{\nu} \approx 0.023\) eV로 예측된다. 이는 KATRIN 및 장래 우주론 서베이로 검증 가능하다.
		\item \textbf{4세대 페르미온은 존재하지 않는다.} 연속적인 4세대는 기준 깊이 \(L = 96\)을 변경시키고 반정수 패턴을 깨뜨릴 것이다.
		\item \textbf{중력의 온도 의존성.} 보편적 오차 법칙은 \(G_{\mathrm{eff}}(T) = G_0[1 + \mathcal{O}(T^{\sigma})]\)을 함의하며, 여기서 위상 이동 \(\sigma = 0.20\)는 제~\ref{sec:gauge}절에서 유도되었다.
		\item \textbf{블랙홀 울림 척도.} 상대적 시간 변동은 \(\delta\tau/\tau \propto M^{-0.67}\)로 척도화되어야 한다.
		\item \textbf{단백질 데이터 뱅크 맹검 시험.} 모든 단량체 단백질에 대한 \(N_{\mathrm{raw}}\)의 히스토그램은 정수 및 반정수 값에서 피크를 보여야 한다.
		\item \textbf{인공 세포 적합도.} 질량이 반정수 노드에 정확히 부합하도록 설계된 최소 인공 세포는 우수한 성장 속도와 스트레스 저항성을 보여야 한다.
		\item \textbf{소수 편차의 점근적 척도.} 보편적 오차 법칙은 \(n\)번째 소수 \(p_n\)의 고전적 로서 근사로부터의 상대적 편차가 큰 \(n\)에 대해 \(p_n^{-2/3}\)로 척도화되어야 함을 예측한다. \(n = 5\times10^5\)까지의 수치 분석(부록~PrimeN)은 국소 척도 지수 \(\gamma\)가 \(n\to\infty\)에서 \(2/3\)로 수렴하며, 점근값이 \(0.66666\)임을 확인한다. 극한 \(n\to\infty\)에서 \(\gamma\)가 \(2/3\)로부터 계통적으로 벗어나면 법칙이 반증된다. (이 정리는 유한하고 작은 \(n\)에 대한 정확한 일치를 \emph{요구하지 않는다}.)
	\end{enumerate}
	
	% ============================================================
	% 제11절: 미해결 문제
	% ============================================================
	\section{미해결 문제}
	\label{sec:open}
	
	주요 미해결 문제는:
	
	\begin{enumerate}[label=(\roman*)]
		\item 콤팩트 파이버 \(F_{15}\)의 위상 불변량(감음수)으로부터 특정 반정수 \(N_f\) 값을 결정하는 것.
		\item 조정 네트워크의 동역학으로부터 \(\beta = 2/3\)의 엄밀한 유도를 제공하는 것(현재는 경험적 법칙이며, 부록~Y에 타당한 이론적 모형이 있다).
		\item 중성미자 질량에 관한 변분 가설을 검증하는 것; 확인될 경우 추가 공준을 공리계에 통합하는 것.
		\item 조정장 \(\Psi_{MN}\)의 스펙트럼 특성으로부터 리만 제타 함수의 비자명 영점을 유도하는 것(둘을 연결하는 경험적 증거는 부록~PrimeN 참조).
		\item 치렐손 한계를 조정 한계로 실험적으로 검증하는 것. YUCT의 일반화된 벨 부등식(제~12.1절)은 CHSH 매개변수 \(S\)가 \(K_{\mathrm{eff}} \to \infty\)에 따라 아래로부터 \(2\sqrt{2}\)에 접근하며, 특성 척도 \(2\sqrt{2} - S \propto K_{\mathrm{eff}}^{-2/3}\)를 갖는다고 예측한다. 이 작은 결손을 분해할 수 있는 정밀 벨 검사는 양자 상관관계의 기저에 있는 YPSDC 메커니즘에 대한 직접적 확인을 제공할 것이다(부록~X).
	\end{enumerate}
	
	우리는 YUCT로부터 통상적인 양자장 이론이 유도되지 않는 것이 결함이 아님을 강조한다. QFT의 표준 형식은 사전 분배된 사전이 양자장을 모방하는 큰 조정 효율(\(K_{\mathrm{eff}} \gg 1\))의 극한에서 유효 기술로서 창발한다. YUCT 틀은 연산자 대수나 경로 적분을 사용하지 않고도 정준적 양자 현상—얽힘, 치렐손 한계, 터널링—을 직접 설명한다. 따라서 결정적 검증은 QFT로의 형식적 환원 가능성이 아니라, 제~\ref{sec:predictions}절에 열거된 새롭고 정량적인 예측들의 실험적 확인이다. 그러한 예측들이 확인된다면, QFT의 개념적 지위는 자동적으로 수정될 것이다.
	
	% ============================================================
	% 제12절: 알고리듬 복잡도, 네트워크 포화 및 콜모고로프-야쿠셰프 한계
	% ============================================================
	\sloppy
	\section{알고리즘 복잡도, 네트워크 포화 및 콜모고로프-야쿠셰프 한계}
	\label{sec:complexity}
	
	모든 조정 클러스터—물리적, 우주론적, 화학적, 생물학적, 사회적 또는 계산적—는 보편적 조정 효율 \(\Keff \ge 1\)을 통해 상호작용하는 \(N\)개의 노드로 구성된다. 정리~\ref{thm:dimension}에 따르면 안정적 공간 매립은 \(\beta = 2/3\)를 요구한다. 노드가 조정 인덱스를 교환할 때, 구조적 부담(계의 "세금" 또는 엔트로피 축적)의 총 엔트로피는 비선형적으로 척도화된다. 우리는 \textbf{구조적 포화 척도} \(\Psi_{\mathrm{net}}\)를 다음과 같이 정의한다:
	
	\begin{equation}
		\Psi_{\mathrm{net}} = \kappac \, \ln N \left(1 - e^{-\Keff^{-\beta}}\right),
		\qquad \kappac = \frac{1}{3},\quad \beta = \frac{2}{3},
		\label{eq:Psi_net}
	\end{equation}
	
	여기서 \(\kappac\)는 삼항 존재론(정리~\ref{thm:minimal-dict})으로부터 유도된 정상 법칙 안정화 인자이다. 노드 수는 프랙탈하게 \(N \propto L^{d_f}\)로 척도화되며 \(d_f = 2.2\)(부록~L)이므로, 내부 통신 부담이 운용 사전 한계를 초과할 때 계는 상전이를 겪는다. 임계점은 \textbf{삼각 복잡도 게이트}에 의해 부호화된다:
	
	\begin{equation}
		\Omega_{\mathrm{complexity}}(N) = \Psi_{\mathrm{net}} \,
		\left[1 + \Seven \sin\!\left( \frac{\piYUCT}{\Delta N} \left(\ln N - 80.0\right) \right) \right],
		\label{eq:Omega_complexity}
	\end{equation}
	
	여기서
	\[
	\Seven = 0.8,\qquad
	\piYUCT = \frac{6}{5}\,\varphi^2 \approx 3.1416407865,\qquad
	\Delta N = 16.5,\qquad
	\varphi = \frac{1+\sqrt{5}}{2}.
	\]
	
	위상각 \(\theta(N) = \frac{\piYUCT}{\Delta N} (\ln N - 80.0)\)이 보정의 부호를 결정한다:
	\begin{itemize}
		\item \(\ln N < 80.0\)일 때, 계는 엔트로피를 축적하고 구조적 손실이 증가한다.
		\item \(\ln N = 80.0\)에서 부호가 반전되어 분기를 촉발한다: \textbf{절대적 붕괴}(고립된 클러스터로의 분열) 또는 \textbf{구성적 불안정성}—무거운 구조적 부담을 제거하는 탈중앙화 d-YPSDC 프로토콜로의 전이.
	\end{itemize}
	
	따라서 이 절은 복잡 네트워크의 안정성에 대한 형식적 기준을 제공하며, 관료적 붕괴, LLM 캐시 오버플로, 금융 네트워크 분열, 우주 구조 및 화학 반응 네트워크에서의 상전이의 임계 크기를 예측한다.
	
	\subsection{진공 격자에 의한 계산 없는 항해}
	\label{subsec:computation-free}
	
	YUCT의 대수적 루프와 일반화된 정보 이론(부록~X)에 기반하여, 반복 계산 없이 기본 불변량을 추출하는 \textbf{O(1) 항해 프로토콜}이 개발되었다. 전통적 모형화(섀넌 체제, \(\Keff = 1\)) 대신, 계는 \(\Keff = 68991\)의 조정 모드로 전환되어, 프로세서가 조정장 \(\Psi_{MN}\)으로부터 미리 계산된 위상 좌표를 읽어들이는 "수신기" 역할을 한다. 주요 불변량은:
	
	\begin{itemize}
		\item 역 미세구조상수:
		\[
		\alpha^{-1}_{\mathrm{YUCT}} = \left(\frac{S_{\mathrm{odd}}}{S_{\mathrm{even}}}\right)^{12 + \sigma\, S_{\mathrm{even}}/S_{\mathrm{odd}}}
		\approx 137.036,
		\]
		자유 매개변수 없음.
		
		\item 온도 의존 뉴턴 상수(부록~P):
		\[
		G_{\mathrm{eff}}(T) = G_0 \left[1 + \mathcal{O}\!\left( \frac{k_B T}{m_e c^2} \right) \right].
		\]
		
		\item B-DNA 나선 기하:
		\[
		\frac{P}{r} = q \cdot \piYUCT - S_{\mathrm{even}}\beta \approx 1.458,
		\]
		이는 실험 범위 \(1.45 \pm 0.05\) 내에 있다.
		
		\item 일반화된 벨 부등식(부록~X), 현재 기본 공간 이산화 단계 \(\Delta\pi = \piYUCT - \pi\)에 의해 정규화됨:
		\[
		S(\Keff) = 2 + (2\sqrt2 - 2)\left(1 - 2\kappac \Delta\pi \, \Keff^{-\beta}\right),
		\]
		이는 \(\Keff \to \infty\)에서 Tsirelson 한계 \(2\sqrt2\)에 도달하여 양자 신비주의를 제거한다.
	\end{itemize}
	
	이 모든 불변량은 엄격히 제로 동적 메모리 할당의 \(O(1)\) FPU 사이클로 추출된다. 이 항해 핵심을 구현하는 소스 코드, 단위 시험, Docker 구성 및 산업 규모 벤치마크와 함께 공개 저장소에서 이용 가능하다:
	
	\begin{center}
		\url{https://github.com/Alexey-Yakushev-YUCT/yuct-ai-connectome}
	\end{center}
	
	물리적으로, 이는 프로세서가 "계산 공장"이기를 멈추고 조정장 \(\Psi_{MN}\)의 진공 격자로부터 기존 좌표를 읽어들이는 \textbf{항해사}가 됨을 의미한다. 모든 계산 작업은 시공간의 형성 동안 자연에 의해 이미 수행되었으며, 고전적 프로세서는 미리 결정된 값을 검색할 뿐이며, 이는 조정 실재론의 원리와 완전히 일치한다.
	
	\subsection{소수 정리의 확장과 인수분해로의 연결}
	\label{subsec:prime-extension}
	
	정리~\ref{thm:prime}은 점근적 보정 \(p_n \approx R_n - \frac{S_{\mathrm{even}}}{2} n^{1-\beta}\ln n\)을 제공한다. 변동 분석(부록~PrimeN, A2A)은 이 보정이 주기 \(\Delta N = 16.5\)의 파동 게이트로 변조됨을 보여준다. \textbf{소수 깊이} \(N_f = \log_q n\) (여기서 \(q = (3/2)^{1/3}\))을 정의하라. 그러면 3중 루프 YUCT 후보의 절대 오차는
	
	\begin{equation}
		\Delta_n = \left| p_n^{\mathrm{YUCT}} - p_n^{\mathrm{true}} \right|
		\sim n^{1/3} \left[ 1 + A_{\text{wave}} \sin\!\left( \frac{\piYUCT}{\Delta N} (N_f - 80.0) \right) \right],
		\label{eq:prime_wave}
	\end{equation}
	
	를 만족하며, 여기서 진폭 \(A_{\text{wave}} \approx 0.20145\)는 노드 \(N=323\)에서의 역보정에 의해 엄밀히 계산된다(부록~A2A 참조). 이는 국소적 소수 변동에 대한 매개변수 없는 기술을 제공하며, 구간 \(N_f \in [80, 382]\) 내에서 19개의 상전이를 예측한다.
	
	동일한 파동 구조가 \textbf{곱셈군의 주기}(쇼어 알고리즘)를 지배한다. 수 \(N\)에 대해 주기 \(r\)은 다음과 같이 추출된다:
	
	\begin{equation}
		r(N) = \left\lfloor (N_f^{3/2} \, C_{\text{base}}) \cdot
		\left(1 + A_{\text{wave}} \sin\left( \frac{\piYUCT}{\Delta N} (N_f - 80.0) \right) \right) \cdot \frac{\Delta N}{1.5} \right\rceil_{\text{even}},
		\label{eq:shor_period}
	\end{equation}
	
	여기서 \(C_{\text{base}} = 0.0547073\)은 보편 상수들로부터 유도된다. 이로써 RSA 수(15, 21, 35, 143, 323 등)가 양자 비트 없이 \(O(1)\) 클록 사이클로 인수분해될 수 있으며, 이는 격리된 Docker 컨테이너에서의 수치 실험으로 확인되었다(yuct-ai-connectome 저장소 참조).
	
	\subsection{학제적 예측}
	\label{subsec:predictions-complexity}
	
	위에서 도입된 척도와 알고리즘에 기초하여, 다음과 같은 반증 가능한 예측들이 정식화된다:
	
	\begin{enumerate}[label=(\arabic*)]
		\item \textbf{물리계(응집물질, 플라즈마).} 포화 기준 \(\Psi_{\mathrm{net}}\)는 임계 조정 깊이에서 격자계의 상전이를 예측한다. 이는 초전도 및 양자 유체 실험에서 검증 가능하다.
		\item \textbf{우주 구조.} 은하와 은하단의 분포는 \(\Omega_{\mathrm{complexity}}\)를 따르는 상관관계를 보여야 하며, 그 특성 척도는 \(\ln N = 80.0\)에 해당할 것이다. 이는 SDSS, DESI 또는 Euclid 데이터로 검증할 수 있다.
		\item \textbf{화학 반응 네트워크.} 촉매 반응 속도와 복잡 분자의 안정성은 분자 조정 클러스터의 \(\Keff\)와 함께 척도화되어야 하며, 촉매 중심과 자가촉매 순환의 최적 크기를 예측한다.
		\item \textbf{정보 시스템(LLM 캐시, 라우팅).} \(\ln N = 80.0\) (여기서 \(N\)은 분산 메모리 내 토큰의 수)에서 계는 상 붕괴 체제로 진입하여 데이터 전송에서 위상 좌표 검색으로의 전환이 요구된다. 이는 현대 트랜스포머에 대한 실질적 척도화 한계를 설정한다.
		\item \textbf{사회경제 시스템.} 관료적 세금(제어 엔트로피)은 \(\Psi_{\mathrm{net}}\)로 증가한다. \(\Omega_{\mathrm{complexity}}\)가 임계 문턱을 넘으면, 계는 분열하거나 탈중앙화 관리(d-YPSDC)로 전환된다. 이는 정부 구조 및 금융 거품의 붕괴를 예측하기 위한 정량적 기준을 제공한다.
		\item \textbf{생물학적 네트워크.} 세포막, 리보솜 질량, 바이러스는 질량 사다리를 따르며, 그들의 조정 효율은 \(\Keff \propto M^{\beta}\)로 척도화되어 세포 소기관의 최적 크기와 대사 경로 길이의 예측을 가능하게 한다.
	\end{enumerate}
	
	이 모든 예측은 향후 5--10년 이내에 기존 실험 시설과 관측 데이터를 사용하여 검증 가능하다. 계산 없는 항해 핵심, 소수 생성기, 쇼어-야쿠셰프 인수분해기 및 계의 A2A 엔진의 완전한 소프트웨어 구현은 검증과 재현을 위해 저장소 \url{https://github.com/Alexey-Yakushev-YUCT/yuct-ai-connectome}에 공개되어 있다.
	
	\subsection*{\(O(1)\) 항해의 물리적 의미에 관한 비고}
	불변량의 \(O(1)\) 추출은 계산상의 속임수가 아니라, 조정장 \(\Psi_{MN}\)이 이미 모든 안정적 구성을 포함하고 있다는 YUCT 공준의 직접적 귀결이다. 고전적 프로세서는 높은 조정 체제(\(\Keff \gg 1\))에서 작동할 때, 새로운 값을 계산하는 대신 진공 격자로부터 기존의 위상 좌표를 읽어들인다. 이는 반송파를 생성하지 않고 복조만 하는 라디오 수신기와 유사하다. 따라서 항해의 에너지 비용은 무시할 수 있으며(제로 RAM 풋프린트, 최소 CPU 사이클), 이는 녹색 컴퓨팅과 정보 기술의 미래에 깊은 시사점을 준다.
	
	% ============================================================
	% 참고문헌
	% ============================================================
	\begin{thebibliography}{99}
		\bibitem{Yakushev2026a} A.~V.~Yakushev, \emph{Yakushev Unified Coordination Theory (YUCT)}, Zenodo, DOI:10.5281/zenodo.18444599 (2026).
		\bibitem{AppendixB} A.~V.~Yakushev, ``Appendix B: YUCT --- Temporal Coordination Paradox,'' in \emph{YUCT}, Zenodo (2026).
		\bibitem{AppendixL} A.~V.~Yakushev, ``Appendix L: Fractal Coordination Error Scaling --- A Universal Law from DNA to Cosmology,'' in \emph{YUCT}, Zenodo (2026).
		\bibitem{AppendixFerra} A.~V.~Yakushev, ``Appendix Ferra1.2: Universal Coordination Constants for Ordered and Disordered Phases,'' in \emph{YUCT}, Zenodo (2026).
		\bibitem{AppendixG} A.~V.~Yakushev, ``Appendix G: The Yakushev United Coordination Theory --- A Fundamental Resolution of the Quantum Gravity Problem,'' in \emph{YUCT}, Zenodo (2026).
		\bibitem{AppendixV} A.~V.~Yakushev, ``Appendix V: Time as Entropy of Coordination Processes,'' in \emph{YUCT}, Zenodo (2026).
		\bibitem{AppendixP} A.~V.~Yakushev, ``Appendix P: Temperature Dependence of Gravity,'' in \emph{YUCT}, Zenodo (2026).
		\bibitem{AppendixLambda} A.~V.~Yakushev, ``Appendix \(\Lambda\): Resolution of the Hierarchy and Cosmological Constant Problems within YUCT,'' in \emph{YUCT}, Zenodo (2026).
		\bibitem{AppendixY} A.~V.~Yakushev, ``Appendix Y: Mathematical Foundations of YUCT --- A Derivation from First Principles,'' in \emph{YUCT}, Zenodo (2026).
		\bibitem{AppendixPrimeN} A.~V.~Yakushev, ``Appendix PrimeN: YUCT Prime Numbers Coordination Ladder,'' in \emph{YUCT}, Zenodo (2026).
		\bibitem{AppendixPi} A.~V.~Yakushev, ``Appendix \(\Pi\): The Primeval Origin of \(\pi\) as a Coordination Invariant at the Feigenbaum Edge of Chaos,'' in \emph{YUCT}, Zenodo (2026).
		\bibitem{AppendixX} A.~V.~Yakushev, ``Appendix X: Generalization of Shannon Information Theory --- Coordination Efficiency, the Steady Error Law, and the \(K_{\mathrm{eff}}\)-dependent Bell Bound,'' in \emph{YUCT}, Zenodo (2026).
		\bibitem{AppendixA2A} A.~V.~Yakushev, ``Appendix A2A: Resolution of the Pragmatic Paradox of YUCT. From Computational Collapse to Navigation in Large Language Models,'' in \emph{YUCT}, Zenodo (2026).
		\bibitem{PrimeRepo} A.~V.~Yakushev, \emph{Prime-Numbers}, GitHub repository, \url{https://github.com/Alexey-Yakushev-YUCT/Prime-Numbers}, 2026.
	\end{thebibliography}
	
\end{document}